\chapter{Scripts et Automatisation}
\label{annex:scripts}

Cette annexe regroupe les scripts principaux développés dans le cadre du projet.

\section{Script de bootstrap PKI Vault}
\label{sec:script_pki}

\lstinputlisting[language=bash, caption=bootstrap-pki-api.sh — Initialisation complète de la PKI Vault]{../../pki/vault/scripts/bootstrap-pki-api.sh}

\section{Script d'émission de certificat Mosquitto}

\lstinputlisting[language=bash, caption=issue-mosquitto-cert-api.sh — Émission du certificat du broker]{../../pki/vault/scripts/issue-mosquitto-cert-api.sh}

\section{Script de génération des certificats de dispositifs}

\lstinputlisting[language=bash, caption=generate-device-certs.sh — Génération des certificats IoT]{../../iot/device-sim/generate_device_certs.sh}

\section{Script de test de connectivité}

\lstinputlisting[language=bash, caption=test-connectivity.sh — Vérification E2E]{../../tests/scenarios/test-connectivity.sh}

\section{Script de performance MQTT}

\begin{lstlisting}[language=bash, caption=Exécution du test de performance]
# Mode simulation (sans broker réel)
cd tests/performance
python mqtt_perf_test.py --mode simulate

# Mode réel (avec broker GNS3)
python mqtt_perf_test.py --mode real \
  --broker 192.168.20.10 \
  --port 8883 \
  --cafile ../../pki/certs/ca/root-ca.crt \
  --certfile ../../pki/certs/services/client.crt \
  --keyfile ../../pki/certs/services/client.key
\end{lstlisting}
